% !TeX root = ../../main.tex
\section{関数}
\subsection{関数}

\begin{definitionbox}{\textbf{関数}}
    \label{def:function}
    2つの数量 $x$, $y$ があり,$x$ の値が決まると,その $x$ に対応する $y$ の値が\fitblankbf[]{ただ一つ}に決まるとき,
    $y$ は $x$ の\fitblankbf[]{関数}であるという.
\end{definitionbox}

\begin{definitionbox}{\textbf{変域}}
    \label{def:domain}
    変数の取りうる値の範囲を,その変数の \fitblankbf[]{変域} という.

    特に,$y$ が $x$ の関数であるとき,$x$ の変域をその関数の \fitblankbf[]{定義域} といい,
    定義域内の $x$ の値に対応する $y$ の変域を \fitblankbf[]{値域} という.
\end{definitionbox}

簡単にいうと,$x$ という \textbf{入力} に対して,関数が何らかの \textbf{操作} をし,その結果として $y$ という \textbf{出力} を返すものをいう.

例えば,2つの数量 $x$, $y$ に対して,$y = 2x$ という関係が成り立っているときのことを考える.
この関数に対して,$x = 3$ という入力を与えると,$y = \fitblank[]{6}$ という出力が得られる.

\begin{techniquebox}{\textbf{関数のTech}}
    \label{tech:function-technique}
    Definition \ref{def:function} より,関数において,一般に次のことが成り立つ.
    \begin{enumerate}
        \item 「\textbf{〜のとき}」や「\textbf{〜を通る}」は \fitblankbf[]{代入} の合図
        \item 「\textbf{平行}」は \fitblankbf[]{傾きが等しい} という意味
        \item 「\textbf{交点}」は \fitblankbf[]{連立方程式} を解きなさい! という意味
    \end{enumerate}
\end{techniquebox}

\subsection{関数の種類}

中学数学で学習する関数は,たかだか \fitblankbf[]{4種類} しかない.

\begin{enumerate}[itemsep=0.8em, topsep=0.2em, parsep=0pt, partopsep=0pt]
    \item \fitblankbf[比例]{6}：\ $\bm{y = ax}$
    \item \fitblankbf[反比例]{6}：\ $\bm{y = \frac{a}{x}}$
    \item \fitblankbf[1次関数]{6}：\ $\bm{y = ax + b}$
    \item \fitblankbf[2次関数]{6}：\ $\bm{y = ax^2}$ $^{*}$
\end{enumerate}

\vfill
\noindent\rule{\textwidth}{0.4pt}

\vspace{0.2em}
\noindent
$^{*}$ 厳密には,中学数学における\textbf{二次関数}は,$y = ax^2$（$a \neq 0$）の形の関数,すなわち \textbf{2乗に比例する関数} を指す.

\newpage
